\chapter{Constructing a controlled latent-state estimator}\label{ch:latentestimation}
\section{The estimation problem the product must own}
A planning library generally expects transition, observation and reward
models. A bank has transactions, incomplete measurements and a historical
decision process. Bridging those objects is a substantive product capability.
The estimation module must specify a latent-state model, reconstruct its
event clock, fit the observed-data likelihood, quantify uncertainty and
state the conditions under which changing the action mechanism has a
causal interpretation. Passing a fitted transition matrix to a
POMDP solver leaves most of that work unresolved.

For $K$ states, $M$ signal categories and a finite action set, the model
has $K-1+|\mathcal A|K(K-1)+K(M-1)$ free probabilities before
reward, covariate and logging parameters. Small $K$ makes exact
marginalization possible and is appropriate for a first implementation.
Increasing $K$ creates flexibility, weakly separated states and more local
optima; state count is a model-selection question, not a quality score.
A state name such as ``fragile liquidity'' is an economic hypothesis.
Ordering states by expected delinquency fixes a label convention; it does
not establish a unique financial mechanism.

\section{Forward and backward messages}
Suppress a separate reward emission and define the observed transition kernel
\[
 M_t(s,s')=q_\eta(a_t\mid s,h_t)
          P_\theta(s'\mid s,a_t)B_\theta(o_{t+1}\mid s',a_t).
\]
When logging uses only observed history and its parameters are distinct,
its factor is constant across latent paths and can be omitted from the
objective for environment parameters. This is a joint-likelihood
factorization argument. It does not assert that adaptively chosen actions
are independent of observations, or that the remaining factor is the
ordinary conditional density of observations given the entire future
action sequence.

Let $\alpha_t$ be the normalized filter after the current signal and
preceding observations and actions. Start with
$c_0=\sum_s\mu(s)B(o_0\mid s)$ and
$\alpha_0(s)=\mu(s)B(o_0\mid s)/c_0$. Then
\begin{align}
 \widetilde\alpha_{t+1}(s')&=\sum_s\alpha_t(s)M_t(s,s'),&
 c_{t+1}&=\sum_{s'}\widetilde\alpha_{t+1}(s'),\\
 \alpha_{t+1}(s')&=\widetilde\alpha_{t+1}(s')/c_{t+1},&
 \log L&=\sum_{t=0}^{T}\log c_t.
 \label{eq:scaled-forward}
\end{align}
Consistently scaled backward messages satisfy $\beta_T(s)=1$ and
$\beta_t(s)=\sum_{s'}M_t(s,s')\beta_{t+1}(s')/c_{t+1}$.
The smoothed transition and state probabilities are
\begin{equation}
 \xi_t(s,s')=
 \frac{\alpha_t(s)M_t(s,s')\beta_{t+1}(s')}{c_{t+1}},
 \qquad \gamma_t(s)=\sum_{s'}\xi_t(s,s'),
 \label{eq:controlled-smoothing}
\end{equation}
with $\gamma_T=\alpha_T$. The identities can be checked by direct
enumeration on a short sequence. They extend the expected-count
construction described by \citet[pp.264--266]{rabiner1989} to the
declared controlled kernel. A separate reward emission multiplies $M_t$
when rewards convey additional state information.

The distinction between filtering and smoothing matters operationally.
Filtering uses information available at the decision. Smoothing uses later
observations and actions to revise the historical state estimate. A plot
of smoothed fragility may help diagnose a fitted model, but using that
series as a feature in a retrospective ``prospective'' policy evaluation
would leak future information.

\section{The M step is a constrained count calculation}
The expected complete log likelihood contains
\[
 \sum_s N_s^0\log\mu_s+
 \sum_{a,s,s'}N^P_{ass'}\log P_a(s,s')+
 \sum_{s,o}N^B_{so}\log B_s(o)+
 \sum_{s,a}N^q_{sa}\log q_s(a).
\]
The counts sum $\gamma$ and $\xi$ over trajectories with the
appropriate observed action or signal. A Lagrange multiplier shows that
maximizing $\sum_jN_j\log p_j$ subject to $\sum_jp_j=1$ gives
$p_j=N_j/\sum_kN_k$. Row-wise application updates the initial law,
transitions, emissions and, when modeled, logging probabilities.
Structural zeros must follow from the model; small positive pseudocounts
are disclosed regularization.

The E step allocates each trajectory probabilistically among states.
The M step updates the specified likelihood from those expected counts.
A Viterbi path instead chooses a single most probable state sequence.
Treating that path as observed data discards state uncertainty and generally
does not maximize the same likelihood. A convenient visualization must
not silently replace the estimator.

EM need not reach a global maximum. Multiple starts, likelihood increments,
iteration caps, boundary probabilities and state separation belong in
the fit receipt. Direct differentiation of the scaled marginalized
likelihood is a useful challenger. Bayesian inference adds priors and
posterior computation; variational estimation introduces approximation
error. None supplies causal identification through numerical convergence.

\section{Unknown emissions and hidden logging: a fitted study}
The study generates 800 independent trajectories with twelve actions,
two latent states and three signal categories:
\[
 \mu=(.72,.28),\qquad
 P_0=\begin{pmatrix}.93&.07\\.22&.78\end{pmatrix},\quad
 P_1=\begin{pmatrix}.94&.06\\.52&.48\end{pmatrix},\quad
 B=\begin{pmatrix}.78&.18&.04\\.16&.34&.50\end{pmatrix}.
\]
The learner receives only signals and actions. Four starts estimate the
initial law, transitions and emissions; labels are aligned by expected
signal severity before parameter errors are evaluated. The generator's
latent states never enter estimation.

Randomized logging chooses action one with probability $0.4$. In the
observational design its probabilities are $0.15$ and $0.80$ in the two
hidden states. The same observational trajectories are fitted with and
without a state-dependent logging factor. The joint fit estimates that
factor as well. The experiment stipulates the finite model and invariance
of transitions and observations under an action intervention. Those
restrictions make the generated environment causal; the estimation code
does not discover their truth.

\begin{table}[htbp]\centering\small
\begin{tabular}{@{}lrrr@{}}
\toprule
Estimation design & Max $P$ error & Max $B$ error & Policy gain (USD)\\
\midrule
Randomized & 0.026 & 0.054 & 24.90\\
Hidden logging, omitted action law & 0.384 & 0.066 & -22.40\\
Hidden logging, joint action law & 0.045 & 0.007 & 13.61\\
\bottomrule
\end{tabular}

\caption{Unknown-emission estimation. The policy contrast is assistance
after the most severe current signal versus no assistance, over eight
decisions. The true gain is USD\,22.34.}
\label{tab:hmm-estimation}
\end{table}

The omitted-action model reverses the sign of the policy contrast in this
sample. Historical assistance reveals fragility; ignoring that selection
can interpret a selected post-assistance distribution as a harmful action
effect. The joint model corrects this particular mechanism, but its
single-sample value remains materially different from truth. Raw objective
values from the two fitting routes are not directly comparable because
the joint objective includes the action density.

\begin{table}[htbp]\centering\small
\begin{tabular}{@{}lrrr@{}}
\toprule
Action model & Gain bias (USD) & Gain RMSE & Transition RMSE\\
\midrule
Omitted & -38.95 & 39.08 & 0.203\\
Joint & 3.41 & 21.90 & 0.051\\
\bottomrule
\end{tabular}

\caption{Thirty independently generated observational datasets of 500
trajectories, with both fitting routes applied to each dataset. Correct
logging specification improves parameter recovery but leaves substantial
policy-value variability.}
\label{tab:hmm-repeated}
\end{table}

The omitted model's mean policy-gain error is approximately
USD\,$-38.95$, with Monte Carlo standard error USD\,0.59. The joint
model's mean error is USD\,3.41 with Monte Carlo standard error USD\,4.02,
and RMSE USD\,21.90. These results expose a failure and a conditional
remedy. They do not show that every controlled hidden-state model is
identified or that the remaining uncertainty permits a bank decision.
The runs retain finite iteration budgets. All 30 selected omitted-action
fits and seven of the 30 joint fits hit the 120-iteration repetition cap.
The repeated table therefore compares the implemented finite-budget
procedures, not fully optimized population estimators. Its error contains
optimization as well as sampling and specification effects.

A separate same-data convergence check extends the main observational
fits to a maximum of 2,000 iterations without changing data or starting
seeds. The omitted-action fit meets the stopping criterion at 163 iterations
for its best start and still gives a policy gain near negative USD\,22.40;
the joint fit again gives positive USD\,13.61. Thus the main sign reversal
persists after the stated numerical stopping test. This does not establish
a global optimum or isolate all specification bias in the repeated study.

\section{Uncertainty includes the observation model}
Seventy-nine trajectory bootstrap resamples of the randomized dataset
refit initial probabilities, transitions and emissions, holding the two
policies fixed. The percentile interval for their value difference is
approximately USD\,$[14.61,38.91]$. The bootstrap count gives coarse
tail quantiles; one selected bootstrap fit hits its 120-iteration cap.
This is a model-conditional interval for independent
trajectories, not a coverage study or an interval for unrestricted hidden
confounding in bank records.

Additional uncertainty arises from the state specification, intervention
invariance and policy selection. A bootstrap holding those choices fixed
does not address them. A fitted three-state model may describe historical
observations as well as a two-state model and value an untested policy
differently. An assistance action may alter reporting intensity as well
as liquidity. A policy selected from hundreds of candidates may exploit
noise. Each requires a separate comparison, sensitivity analysis or outer
evaluation design.

The product should show state separation, transition support, likelihood
profiles or posterior dependence, predictive residuals by regime, and
financial decision sensitivity. A large parameter interval need not imply
a large value interval, and a small parameter error can reverse a close
policy comparison. The decision determines which uncertainty matters.

\section{Planning with a distribution over states}
For a correctly specified environment and an actor restricted to recorded
history, let $b$ be its state belief, $p(o'\mid b,a)$ the next-signal
law and $\tau(b,a,o')$ its Bayes update. With feasible actions
$\mathcal A(b)$ and expected reward $r(b,a)$,
\begin{equation}
 V_t(b)=\max_{a\in\mathcal A(b)}\left[
 r(b,a)+\gamma\sum_{o'}p(o'\mid b,a)
       V_{t+1}\{\tau(b,a,o')\}\right].
 \label{eq:belief-bellman}
\end{equation}
Observed eligibility must be part of the information state if it changes
the action set. The continuation averages future state uncertainty and
future information. A myopic threshold neglects both continuation rewards
and any value of information \citep{kaelbling1998}.

The study values next-signal contribution at $(18,0,-90)$ dollars and
charges eight dollars for assistance. Reward is a deterministic function
of the next signal and action, so no separate reward information is
omitted from the filter. A 501-point belief grid solves eight-decision
control. Doubling resolution changes initial value by less than
USD\,0.0001. Under the supplied true model, values are approximately
$-21.92$ for no assistance, $0.42$ for assistance after a severe signal,
and $5.30$ for the grid policy. These numerical comparisons do not show
that the extra optimized value is statistically resolved by the fitted data.

\begin{figure}[htbp]\centering
\includegraphics[width=.9\textwidth]{derivations/pomdp_value.pdf}
\caption{Eight-decision value and initial action on a two-state belief grid,
conditional on the supplied causal environment.}
\label{fig:pomdp-plan}
\end{figure}

For bounded rewards and $\gamma<1$, the infinite-horizon Bellman
operator is a contraction of modulus $\gamma$: expectation changes by
at most the supremum change in continuation; discounting multiplies the
bound by $\gamma$; taking a maximum preserves it. This establishes a
unique fixed point for the supplied belief-MDP. It is a control result,
not an identification theorem for observational banking data.

\section{Partial observation is not always hidden confounding}
Under randomized or recorded-history assignment, partial observation can
make modeling difficult without confounding the action. History-based
importance evaluation remains identified under the sequential assumptions,
although its variance may be prohibitive. A latent model can impose
structure that improves precision, with a corresponding specification risk.

If historical staff used unrecorded information predictive of future
outcomes, an accurate $\Prob(A_t\mid H_t)$ is generally insufficient.
It describes selected historical actions without reconstructing hidden
information under a new policy. \citet{tennenholtz2020,bennettprox2023}
study this problem using specified information structures and informative
proxy or bridge conditions. Their results do not turn every fitted HMM
into a causal environment.

The empirical questions are about measurement: which private information
did staff use; which independently recorded proxies have defensible roles;
could limited randomization identify a missing transition; could a narrower
policy be evaluated with existing support? These questions can be more
valuable than increasing the representation dimension.

\subsection{A finite bridge exposes the inverse problem}
At fixed $(a,x)$ in Equation~\eqref{eq:bridge}, let
$M_{zw}=\Prob(W=w\mid Z=z,A=a,X=x)$ and
$m_z=\E[Y\mid Z=z,A=a,X=x]$. A discrete bridge solves
$M\boldsymbol h=\boldsymbol m$. Full column rank gives at most one
bridge in this class. A small minimum singular value lets modest noise
in $\boldsymbol m$ produce large changes in $\boldsymbol h$.
Regularization stabilizes the calculation while introducing bias.

This explains why two predictive but almost redundant proxies may be weak
for identification. It also limits a numerical diagnostic: empirical full
rank does not establish the substantive proxy exclusions or population
completeness. Sequential bridge methods additionally preserve temporal
roles and the target actor's information. The time-independent-sampling
identification construction and proximal estimating strategy discussed in
\citet{bennettprox2023} are distinct approaches, with different computational
and regularity requirements.

\section{Customer behavior and the bank's objective}
A bank POMDP is not automatically a model of customer utility. Distinguish
bank offer $b$ from household choice $a$. A declared dynamic logit
construction has
$v(a,s;b)=u(a,s;b)+\beta\E[\overline V(s')\mid s,a,b]$ and,
under centered independent Type-I extreme-value shocks of fixed scale
$\tau$, $\overline V(s)=\tau\log\sum_a\exp\{v(a,s)/\tau\}$.
Choice probabilities are proportional to $\exp\{v(a,s)/\tau\}$.
Utility normalization, discounting, information and transition restrictions
must be identified; a softmax-looking policy does not supply them.

Given a supported customer response law, the bank's induced objects are
\[
 K_b(s,s')=\sum_aP_\theta(a\mid s,b)F_\theta(s'\mid s,a,b),
 \qquad r_b(s)=\sum_aP_\theta(a\mid s,b)m(s,b,a).
\]
Household utility $u$ and bank contribution $m$ have different owners
and need not be proportional. Mortgage demand and lender behavior in
\citet{benetton2021}, and refinancing mechanisms in
\citet{agarwal2020,beraja2019}, motivate this separation. The synthetic
POMDP's supplied reward schedule does not estimate household utility.

Composition therefore needs three contracts: the response module declares
supported offer interventions; finance declares ownership, discounting
and terminal claims; control declares the bank's information and feasible
actions. A fitted incumbent-choice probability alone does not support an
unmodeled marketing or pricing intervention.

\section{A concrete initial module}
The initial product should support finite states, categorical and simple
parametric emissions, observed covariates, variable trajectory lengths,
explicit censoring and a declared logging-law family. Results contain
likelihood diagnostics, parameter and state uncertainty, transition support,
the observation clock and a separate causal interpretation. Prospective
trajectory views use filtering; retrospective views label smoothing.

Small models compare EM, direct optimization and exact path enumeration.
Planning consumes an explicitly accepted environment and preserves the
actor's information. Evaluation fixes a policy or runs an outer training
and evaluation protocol. Continuous states, particle likelihoods and
amortized inference have separate approximation and identification work.
BayesFilter can supply useful probabilistic infrastructure, but the causal
action model, intervention evidence and temporal semantics remain owned by
this product. No current bank or sibling integration is claimed.
